% ============================================================
% fig_transformer_completo.tex — Arquitectura Transformer completa
% Uso: % ============================================================
% fig_transformer_completo.tex — Arquitectura Transformer completa
% Uso: % ============================================================
% fig_transformer_completo.tex — Arquitectura Transformer completa
% Uso: % ============================================================
% fig_transformer_completo.tex — Arquitectura Transformer completa
% Uso: \input{figures/fig_transformer_completo}
% ============================================================
\begin{figure}[H]
\centering
\begin{tikzpicture}[
    blk/.style={draw, rounded corners, minimum width=3.0cm,
                minimum height=0.62cm, font=\footnotesize, align=center},
    enc/.style={blk, fill=blue!10, draw=blue!40},
    dec/.style={blk, fill=green!10, draw=green!50!black},
    norm/.style={blk, fill=orange!10, draw=orange!50,
                 minimum height=0.5cm},
    emb/.style={blk, fill=purple!10, draw=purple!40},
    out/.style={blk, fill=red!8, draw=red!40},
    add/.style={draw, circle, fill=yellow!20, draw=yellow!60!black,
                minimum size=0.42cm, font=\scriptsize\bfseries},
    arrow/.style={-{Stealth}, thick},
    skip/.style={-{Stealth}, thick, gray!50, dashed},
    lbl/.style={font=\tiny, gray}
]

%═══════════════════════════════════════════════════════════
% ENCODER (izquierda)
%═══════════════════════════════════════════════════════════
\node[font=\small\bfseries, blue!60] at (-3.2, 10.8) {ENCODER};

% Entrada
\node[lbl] (src) at (-3.2, -0.4) {secuencia de entrada};
\node[emb] (src_emb) at (-3.2, 0.55) {Input Embedding};
\node[emb] (src_pe)  at (-3.2, 1.45) {Positional Encoding};

% Bloque encoder (representado como caja grande)
\draw[blue!30, rounded corners, thick, fill=blue!3]
    (-4.9, 2.1) rectangle (-1.5, 7.9);
\node[lbl] at (-4.55, 7.65) {$\times N$};

\node[enc]  (mha_e)  at (-3.2, 2.75) {Multi-Head\\Self-Attention};
\node[add]  (add1e)  at (-3.2, 3.75) {$+$};
\node[norm] (ln1e)   at (-3.2, 4.4)  {Add \& Norm};
\node[enc]  (ffn_e)  at (-3.2, 5.3)  {Feed-Forward};
\node[add]  (add2e)  at (-3.2, 6.3)  {$+$};
\node[norm] (ln2e)   at (-3.2, 6.95) {Add \& Norm};

% Salida encoder
\node[lbl] (enc_out) at (-3.2, 8.3) {$K, V$ al decoder};

% Flechas encoder
\draw[arrow] (src)     -- (src_emb);
\draw[arrow] (src_emb) -- (src_pe);
\draw[arrow] (src_pe)  -- (mha_e);
\draw[arrow] (mha_e)   -- (add1e);
\draw[arrow] (add1e)   -- (ln1e);
\draw[arrow] (ln1e)    -- (ffn_e);
\draw[arrow] (ffn_e)   -- (add2e);
\draw[arrow] (add2e)   -- (ln2e);
\draw[arrow] (ln2e)    -- (enc_out);

% Residuales encoder
\draw[skip] (src_pe.west) -- ++(-0.7,0) -- ++(0,3.75) -- (add1e.west);
\draw[skip] (ln1e.west)   -- ++(-0.5,0) -- ++(0,2.0)  -- (add2e.west);

%═══════════════════════════════════════════════════════════
% DECODER (derecha)
%═══════════════════════════════════════════════════════════
\node[font=\small\bfseries, green!60!black] at (3.2, 10.8) {DECODER};

% Entrada
\node[lbl] (tgt) at (3.2, -0.4) {secuencia de salida (desplazada)};
\node[emb] (tgt_emb) at (3.2, 0.55) {Output Embedding};
\node[emb] (tgt_pe)  at (3.2, 1.45) {Positional Encoding};

% Bloque decoder
\draw[green!30, rounded corners, thick, fill=green!3]
    (1.5, 2.1) rectangle (4.9, 9.1);
\node[lbl] at (4.55, 8.85) {$\times N$};

\node[dec]  (mha_d1) at (3.2, 2.75) {Masked Multi-Head\\Self-Attention};
\node[add]  (add1d)  at (3.2, 3.75) {$+$};
\node[norm] (ln1d)   at (3.2, 4.4)  {Add \& Norm};
\node[dec]  (mha_d2) at (3.2, 5.3)  {Multi-Head\\Cross-Attention};
\node[add]  (add2d)  at (3.2, 6.3)  {$+$};
\node[norm] (ln2d)   at (3.2, 6.95) {Add \& Norm};
\node[dec]  (ffn_d)  at (3.2, 7.85) {Feed-Forward};
\node[add]  (add3d)  at (3.2, 8.75) {$+$};
% (ln3 dentro del bloque, salida)

% Flechas decoder
\draw[arrow] (tgt)     -- (tgt_emb);
\draw[arrow] (tgt_emb) -- (tgt_pe);
\draw[arrow] (tgt_pe)  -- (mha_d1);
\draw[arrow] (mha_d1)  -- (add1d);
\draw[arrow] (add1d)   -- (ln1d);
\draw[arrow] (ln1d)    -- (mha_d2);
\draw[arrow] (mha_d2)  -- (add2d);
\draw[arrow] (add2d)   -- (ln2d);
\draw[arrow] (ln2d)    -- (ffn_d);
\draw[arrow] (ffn_d)   -- (add3d);

% Residuales decoder
\draw[skip] (tgt_pe.east) -- ++(0.7,0) -- ++(0,3.75) -- (add1d.east);
\draw[skip] (ln1d.east)   -- ++(0.5,0) -- ++(0,2.0)  -- (add2d.east);
\draw[skip] (ln2d.east)   -- ++(0.5,0) -- ++(0,2.0)  -- (add3d.east);

% ── Conexión encoder → decoder (cross-attention) ─────────────
\draw[arrow, red!60, very thick]
    (enc_out.east) -- ++(0.5,0) -- ++(0,-3.0) -- (mha_d2.west);
\node[font=\tiny, red!60, align=center] at (0, 6.8)
    {$K$, $V$ del\\encoder};

% ── Salida final ──────────────────────────────────────────────
\node[out]  (linear)  at (3.2, 9.8)  {Linear};
\node[out]  (softmax) at (3.2, 10.5) {Softmax};
\node[lbl]           at (3.2, 11.1)  {probabilidades de salida};

\draw[arrow] (add3d)  -- (linear);
\draw[arrow] (linear)  -- (softmax);

% ── Nota máscara ─────────────────────────────────────────────
\node[draw=gray!30, fill=gray!5, rounded corners, font=\tiny,
      inner sep=4pt, align=left] at (6.8, 2.75)
    {\textbf{Masked:} m\'ascara causal\\impide ver tokens futuros.\\Garantiza generaci\'on\\autoregresiva.};
\draw[-{Stealth}, gray!40, thin] (4.9, 2.75) -- (5.5, 2.75);

\end{tikzpicture}
\caption{Arquitectura completa del Transformer original~\cite{vaswani2017attention}.
El encoder (izquierda) procesa la secuencia de entrada en paralelo y produce
representaciones contextualizadas. El decoder (derecha) genera la secuencia
de salida token a token, utilizando atenci\'on enmascarada sobre su propia
salida parcial y atenci\'on cruzada (\textit{cross-attention}) sobre las
representaciones del encoder. Las flechas rojas indican el flujo de $K$ y $V$
del encoder al decoder.}
\label{fig:transformer_completo}
\end{figure}

% ============================================================
\begin{figure}[H]
\centering
\begin{tikzpicture}[
    blk/.style={draw, rounded corners, minimum width=3.0cm,
                minimum height=0.62cm, font=\footnotesize, align=center},
    enc/.style={blk, fill=blue!10, draw=blue!40},
    dec/.style={blk, fill=green!10, draw=green!50!black},
    norm/.style={blk, fill=orange!10, draw=orange!50,
                 minimum height=0.5cm},
    emb/.style={blk, fill=purple!10, draw=purple!40},
    out/.style={blk, fill=red!8, draw=red!40},
    add/.style={draw, circle, fill=yellow!20, draw=yellow!60!black,
                minimum size=0.42cm, font=\scriptsize\bfseries},
    arrow/.style={-{Stealth}, thick},
    skip/.style={-{Stealth}, thick, gray!50, dashed},
    lbl/.style={font=\tiny, gray}
]

%═══════════════════════════════════════════════════════════
% ENCODER (izquierda)
%═══════════════════════════════════════════════════════════
\node[font=\small\bfseries, blue!60] at (-3.2, 10.8) {ENCODER};

% Entrada
\node[lbl] (src) at (-3.2, -0.4) {secuencia de entrada};
\node[emb] (src_emb) at (-3.2, 0.55) {Input Embedding};
\node[emb] (src_pe)  at (-3.2, 1.45) {Positional Encoding};

% Bloque encoder (representado como caja grande)
\draw[blue!30, rounded corners, thick, fill=blue!3]
    (-4.9, 2.1) rectangle (-1.5, 7.9);
\node[lbl] at (-4.55, 7.65) {$\times N$};

\node[enc]  (mha_e)  at (-3.2, 2.75) {Multi-Head\\Self-Attention};
\node[add]  (add1e)  at (-3.2, 3.75) {$+$};
\node[norm] (ln1e)   at (-3.2, 4.4)  {Add \& Norm};
\node[enc]  (ffn_e)  at (-3.2, 5.3)  {Feed-Forward};
\node[add]  (add2e)  at (-3.2, 6.3)  {$+$};
\node[norm] (ln2e)   at (-3.2, 6.95) {Add \& Norm};

% Salida encoder
\node[lbl] (enc_out) at (-3.2, 8.3) {$K, V$ al decoder};

% Flechas encoder
\draw[arrow] (src)     -- (src_emb);
\draw[arrow] (src_emb) -- (src_pe);
\draw[arrow] (src_pe)  -- (mha_e);
\draw[arrow] (mha_e)   -- (add1e);
\draw[arrow] (add1e)   -- (ln1e);
\draw[arrow] (ln1e)    -- (ffn_e);
\draw[arrow] (ffn_e)   -- (add2e);
\draw[arrow] (add2e)   -- (ln2e);
\draw[arrow] (ln2e)    -- (enc_out);

% Residuales encoder
\draw[skip] (src_pe.west) -- ++(-0.7,0) -- ++(0,3.75) -- (add1e.west);
\draw[skip] (ln1e.west)   -- ++(-0.5,0) -- ++(0,2.0)  -- (add2e.west);

%═══════════════════════════════════════════════════════════
% DECODER (derecha)
%═══════════════════════════════════════════════════════════
\node[font=\small\bfseries, green!60!black] at (3.2, 10.8) {DECODER};

% Entrada
\node[lbl] (tgt) at (3.2, -0.4) {secuencia de salida (desplazada)};
\node[emb] (tgt_emb) at (3.2, 0.55) {Output Embedding};
\node[emb] (tgt_pe)  at (3.2, 1.45) {Positional Encoding};

% Bloque decoder
\draw[green!30, rounded corners, thick, fill=green!3]
    (1.5, 2.1) rectangle (4.9, 9.1);
\node[lbl] at (4.55, 8.85) {$\times N$};

\node[dec]  (mha_d1) at (3.2, 2.75) {Masked Multi-Head\\Self-Attention};
\node[add]  (add1d)  at (3.2, 3.75) {$+$};
\node[norm] (ln1d)   at (3.2, 4.4)  {Add \& Norm};
\node[dec]  (mha_d2) at (3.2, 5.3)  {Multi-Head\\Cross-Attention};
\node[add]  (add2d)  at (3.2, 6.3)  {$+$};
\node[norm] (ln2d)   at (3.2, 6.95) {Add \& Norm};
\node[dec]  (ffn_d)  at (3.2, 7.85) {Feed-Forward};
\node[add]  (add3d)  at (3.2, 8.75) {$+$};
% (ln3 dentro del bloque, salida)

% Flechas decoder
\draw[arrow] (tgt)     -- (tgt_emb);
\draw[arrow] (tgt_emb) -- (tgt_pe);
\draw[arrow] (tgt_pe)  -- (mha_d1);
\draw[arrow] (mha_d1)  -- (add1d);
\draw[arrow] (add1d)   -- (ln1d);
\draw[arrow] (ln1d)    -- (mha_d2);
\draw[arrow] (mha_d2)  -- (add2d);
\draw[arrow] (add2d)   -- (ln2d);
\draw[arrow] (ln2d)    -- (ffn_d);
\draw[arrow] (ffn_d)   -- (add3d);

% Residuales decoder
\draw[skip] (tgt_pe.east) -- ++(0.7,0) -- ++(0,3.75) -- (add1d.east);
\draw[skip] (ln1d.east)   -- ++(0.5,0) -- ++(0,2.0)  -- (add2d.east);
\draw[skip] (ln2d.east)   -- ++(0.5,0) -- ++(0,2.0)  -- (add3d.east);

% ── Conexión encoder → decoder (cross-attention) ─────────────
\draw[arrow, red!60, very thick]
    (enc_out.east) -- ++(0.5,0) -- ++(0,-3.0) -- (mha_d2.west);
\node[font=\tiny, red!60, align=center] at (0, 6.8)
    {$K$, $V$ del\\encoder};

% ── Salida final ──────────────────────────────────────────────
\node[out]  (linear)  at (3.2, 9.8)  {Linear};
\node[out]  (softmax) at (3.2, 10.5) {Softmax};
\node[lbl]           at (3.2, 11.1)  {probabilidades de salida};

\draw[arrow] (add3d)  -- (linear);
\draw[arrow] (linear)  -- (softmax);

% ── Nota máscara ─────────────────────────────────────────────
\node[draw=gray!30, fill=gray!5, rounded corners, font=\tiny,
      inner sep=4pt, align=left] at (6.8, 2.75)
    {\textbf{Masked:} m\'ascara causal\\impide ver tokens futuros.\\Garantiza generaci\'on\\autoregresiva.};
\draw[-{Stealth}, gray!40, thin] (4.9, 2.75) -- (5.5, 2.75);

\end{tikzpicture}
\caption{Arquitectura completa del Transformer original~\cite{vaswani2017attention}.
El encoder (izquierda) procesa la secuencia de entrada en paralelo y produce
representaciones contextualizadas. El decoder (derecha) genera la secuencia
de salida token a token, utilizando atenci\'on enmascarada sobre su propia
salida parcial y atenci\'on cruzada (\textit{cross-attention}) sobre las
representaciones del encoder. Las flechas rojas indican el flujo de $K$ y $V$
del encoder al decoder.}
\label{fig:transformer_completo}
\end{figure}

% ============================================================
\begin{figure}[H]
\centering
\begin{tikzpicture}[
    blk/.style={draw, rounded corners, minimum width=3.0cm,
                minimum height=0.62cm, font=\footnotesize, align=center},
    enc/.style={blk, fill=blue!10, draw=blue!40},
    dec/.style={blk, fill=green!10, draw=green!50!black},
    norm/.style={blk, fill=orange!10, draw=orange!50,
                 minimum height=0.5cm},
    emb/.style={blk, fill=purple!10, draw=purple!40},
    out/.style={blk, fill=red!8, draw=red!40},
    add/.style={draw, circle, fill=yellow!20, draw=yellow!60!black,
                minimum size=0.42cm, font=\scriptsize\bfseries},
    arrow/.style={-{Stealth}, thick},
    skip/.style={-{Stealth}, thick, gray!50, dashed},
    lbl/.style={font=\tiny, gray}
]

%═══════════════════════════════════════════════════════════
% ENCODER (izquierda)
%═══════════════════════════════════════════════════════════
\node[font=\small\bfseries, blue!60] at (-3.2, 10.8) {ENCODER};

% Entrada
\node[lbl] (src) at (-3.2, -0.4) {secuencia de entrada};
\node[emb] (src_emb) at (-3.2, 0.55) {Input Embedding};
\node[emb] (src_pe)  at (-3.2, 1.45) {Positional Encoding};

% Bloque encoder (representado como caja grande)
\draw[blue!30, rounded corners, thick, fill=blue!3]
    (-4.9, 2.1) rectangle (-1.5, 7.9);
\node[lbl] at (-4.55, 7.65) {$\times N$};

\node[enc]  (mha_e)  at (-3.2, 2.75) {Multi-Head\\Self-Attention};
\node[add]  (add1e)  at (-3.2, 3.75) {$+$};
\node[norm] (ln1e)   at (-3.2, 4.4)  {Add \& Norm};
\node[enc]  (ffn_e)  at (-3.2, 5.3)  {Feed-Forward};
\node[add]  (add2e)  at (-3.2, 6.3)  {$+$};
\node[norm] (ln2e)   at (-3.2, 6.95) {Add \& Norm};

% Salida encoder
\node[lbl] (enc_out) at (-3.2, 8.3) {$K, V$ al decoder};

% Flechas encoder
\draw[arrow] (src)     -- (src_emb);
\draw[arrow] (src_emb) -- (src_pe);
\draw[arrow] (src_pe)  -- (mha_e);
\draw[arrow] (mha_e)   -- (add1e);
\draw[arrow] (add1e)   -- (ln1e);
\draw[arrow] (ln1e)    -- (ffn_e);
\draw[arrow] (ffn_e)   -- (add2e);
\draw[arrow] (add2e)   -- (ln2e);
\draw[arrow] (ln2e)    -- (enc_out);

% Residuales encoder
\draw[skip] (src_pe.west) -- ++(-0.7,0) -- ++(0,3.75) -- (add1e.west);
\draw[skip] (ln1e.west)   -- ++(-0.5,0) -- ++(0,2.0)  -- (add2e.west);

%═══════════════════════════════════════════════════════════
% DECODER (derecha)
%═══════════════════════════════════════════════════════════
\node[font=\small\bfseries, green!60!black] at (3.2, 10.8) {DECODER};

% Entrada
\node[lbl] (tgt) at (3.2, -0.4) {secuencia de salida (desplazada)};
\node[emb] (tgt_emb) at (3.2, 0.55) {Output Embedding};
\node[emb] (tgt_pe)  at (3.2, 1.45) {Positional Encoding};

% Bloque decoder
\draw[green!30, rounded corners, thick, fill=green!3]
    (1.5, 2.1) rectangle (4.9, 9.1);
\node[lbl] at (4.55, 8.85) {$\times N$};

\node[dec]  (mha_d1) at (3.2, 2.75) {Masked Multi-Head\\Self-Attention};
\node[add]  (add1d)  at (3.2, 3.75) {$+$};
\node[norm] (ln1d)   at (3.2, 4.4)  {Add \& Norm};
\node[dec]  (mha_d2) at (3.2, 5.3)  {Multi-Head\\Cross-Attention};
\node[add]  (add2d)  at (3.2, 6.3)  {$+$};
\node[norm] (ln2d)   at (3.2, 6.95) {Add \& Norm};
\node[dec]  (ffn_d)  at (3.2, 7.85) {Feed-Forward};
\node[add]  (add3d)  at (3.2, 8.75) {$+$};
% (ln3 dentro del bloque, salida)

% Flechas decoder
\draw[arrow] (tgt)     -- (tgt_emb);
\draw[arrow] (tgt_emb) -- (tgt_pe);
\draw[arrow] (tgt_pe)  -- (mha_d1);
\draw[arrow] (mha_d1)  -- (add1d);
\draw[arrow] (add1d)   -- (ln1d);
\draw[arrow] (ln1d)    -- (mha_d2);
\draw[arrow] (mha_d2)  -- (add2d);
\draw[arrow] (add2d)   -- (ln2d);
\draw[arrow] (ln2d)    -- (ffn_d);
\draw[arrow] (ffn_d)   -- (add3d);

% Residuales decoder
\draw[skip] (tgt_pe.east) -- ++(0.7,0) -- ++(0,3.75) -- (add1d.east);
\draw[skip] (ln1d.east)   -- ++(0.5,0) -- ++(0,2.0)  -- (add2d.east);
\draw[skip] (ln2d.east)   -- ++(0.5,0) -- ++(0,2.0)  -- (add3d.east);

% ── Conexión encoder → decoder (cross-attention) ─────────────
\draw[arrow, red!60, very thick]
    (enc_out.east) -- ++(0.5,0) -- ++(0,-3.0) -- (mha_d2.west);
\node[font=\tiny, red!60, align=center] at (0, 6.8)
    {$K$, $V$ del\\encoder};

% ── Salida final ──────────────────────────────────────────────
\node[out]  (linear)  at (3.2, 9.8)  {Linear};
\node[out]  (softmax) at (3.2, 10.5) {Softmax};
\node[lbl]           at (3.2, 11.1)  {probabilidades de salida};

\draw[arrow] (add3d)  -- (linear);
\draw[arrow] (linear)  -- (softmax);

% ── Nota máscara ─────────────────────────────────────────────
\node[draw=gray!30, fill=gray!5, rounded corners, font=\tiny,
      inner sep=4pt, align=left] at (6.8, 2.75)
    {\textbf{Masked:} m\'ascara causal\\impide ver tokens futuros.\\Garantiza generaci\'on\\autoregresiva.};
\draw[-{Stealth}, gray!40, thin] (4.9, 2.75) -- (5.5, 2.75);

\end{tikzpicture}
\caption{Arquitectura completa del Transformer original~\cite{vaswani2017attention}.
El encoder (izquierda) procesa la secuencia de entrada en paralelo y produce
representaciones contextualizadas. El decoder (derecha) genera la secuencia
de salida token a token, utilizando atenci\'on enmascarada sobre su propia
salida parcial y atenci\'on cruzada (\textit{cross-attention}) sobre las
representaciones del encoder. Las flechas rojas indican el flujo de $K$ y $V$
del encoder al decoder.}
\label{fig:transformer_completo}
\end{figure}

% ============================================================
\begin{figure}[H]
\centering
\begin{tikzpicture}[
    blk/.style={draw, rounded corners, minimum width=3.0cm,
                minimum height=0.62cm, font=\footnotesize, align=center},
    enc/.style={blk, fill=blue!10, draw=blue!40},
    dec/.style={blk, fill=green!10, draw=green!50!black},
    norm/.style={blk, fill=orange!10, draw=orange!50,
                 minimum height=0.5cm},
    emb/.style={blk, fill=purple!10, draw=purple!40},
    out/.style={blk, fill=red!8, draw=red!40},
    add/.style={draw, circle, fill=yellow!20, draw=yellow!60!black,
                minimum size=0.42cm, font=\scriptsize\bfseries},
    arrow/.style={-{Stealth}, thick},
    skip/.style={-{Stealth}, thick, gray!50, dashed},
    lbl/.style={font=\tiny, gray}
]

%═══════════════════════════════════════════════════════════
% ENCODER (izquierda)
%═══════════════════════════════════════════════════════════
\node[font=\small\bfseries, blue!60] at (-3.2, 10.8) {ENCODER};

% Entrada
\node[lbl] (src) at (-3.2, -0.4) {secuencia de entrada};
\node[emb] (src_emb) at (-3.2, 0.55) {Input Embedding};
\node[emb] (src_pe)  at (-3.2, 1.45) {Positional Encoding};

% Bloque encoder (representado como caja grande)
\draw[blue!30, rounded corners, thick, fill=blue!3]
    (-4.9, 2.1) rectangle (-1.5, 7.9);
\node[lbl] at (-4.55, 7.65) {$\times N$};

\node[enc]  (mha_e)  at (-3.2, 2.75) {Multi-Head\\Self-Attention};
\node[add]  (add1e)  at (-3.2, 3.75) {$+$};
\node[norm] (ln1e)   at (-3.2, 4.4)  {Add \& Norm};
\node[enc]  (ffn_e)  at (-3.2, 5.3)  {Feed-Forward};
\node[add]  (add2e)  at (-3.2, 6.3)  {$+$};
\node[norm] (ln2e)   at (-3.2, 6.95) {Add \& Norm};

% Salida encoder
\node[lbl] (enc_out) at (-3.2, 8.3) {$K, V$ al decoder};

% Flechas encoder
\draw[arrow] (src)     -- (src_emb);
\draw[arrow] (src_emb) -- (src_pe);
\draw[arrow] (src_pe)  -- (mha_e);
\draw[arrow] (mha_e)   -- (add1e);
\draw[arrow] (add1e)   -- (ln1e);
\draw[arrow] (ln1e)    -- (ffn_e);
\draw[arrow] (ffn_e)   -- (add2e);
\draw[arrow] (add2e)   -- (ln2e);
\draw[arrow] (ln2e)    -- (enc_out);

% Residuales encoder
\draw[skip] (src_pe.west) -- ++(-0.7,0) -- ++(0,3.75) -- (add1e.west);
\draw[skip] (ln1e.west)   -- ++(-0.5,0) -- ++(0,2.0)  -- (add2e.west);

%═══════════════════════════════════════════════════════════
% DECODER (derecha)
%═══════════════════════════════════════════════════════════
\node[font=\small\bfseries, green!60!black] at (3.2, 10.8) {DECODER};

% Entrada
\node[lbl] (tgt) at (3.2, -0.4) {secuencia de salida (desplazada)};
\node[emb] (tgt_emb) at (3.2, 0.55) {Output Embedding};
\node[emb] (tgt_pe)  at (3.2, 1.45) {Positional Encoding};

% Bloque decoder
\draw[green!30, rounded corners, thick, fill=green!3]
    (1.5, 2.1) rectangle (4.9, 9.1);
\node[lbl] at (4.55, 8.85) {$\times N$};

\node[dec]  (mha_d1) at (3.2, 2.75) {Masked Multi-Head\\Self-Attention};
\node[add]  (add1d)  at (3.2, 3.75) {$+$};
\node[norm] (ln1d)   at (3.2, 4.4)  {Add \& Norm};
\node[dec]  (mha_d2) at (3.2, 5.3)  {Multi-Head\\Cross-Attention};
\node[add]  (add2d)  at (3.2, 6.3)  {$+$};
\node[norm] (ln2d)   at (3.2, 6.95) {Add \& Norm};
\node[dec]  (ffn_d)  at (3.2, 7.85) {Feed-Forward};
\node[add]  (add3d)  at (3.2, 8.75) {$+$};
% (ln3 dentro del bloque, salida)

% Flechas decoder
\draw[arrow] (tgt)     -- (tgt_emb);
\draw[arrow] (tgt_emb) -- (tgt_pe);
\draw[arrow] (tgt_pe)  -- (mha_d1);
\draw[arrow] (mha_d1)  -- (add1d);
\draw[arrow] (add1d)   -- (ln1d);
\draw[arrow] (ln1d)    -- (mha_d2);
\draw[arrow] (mha_d2)  -- (add2d);
\draw[arrow] (add2d)   -- (ln2d);
\draw[arrow] (ln2d)    -- (ffn_d);
\draw[arrow] (ffn_d)   -- (add3d);

% Residuales decoder
\draw[skip] (tgt_pe.east) -- ++(0.7,0) -- ++(0,3.75) -- (add1d.east);
\draw[skip] (ln1d.east)   -- ++(0.5,0) -- ++(0,2.0)  -- (add2d.east);
\draw[skip] (ln2d.east)   -- ++(0.5,0) -- ++(0,2.0)  -- (add3d.east);

% ── Conexión encoder → decoder (cross-attention) ─────────────
\draw[arrow, red!60, very thick]
    (enc_out.east) -- ++(0.5,0) -- ++(0,-3.0) -- (mha_d2.west);
\node[font=\tiny, red!60, align=center] at (0, 6.8)
    {$K$, $V$ del\\encoder};

% ── Salida final ──────────────────────────────────────────────
\node[out]  (linear)  at (3.2, 9.8)  {Linear};
\node[out]  (softmax) at (3.2, 10.5) {Softmax};
\node[lbl]           at (3.2, 11.1)  {probabilidades de salida};

\draw[arrow] (add3d)  -- (linear);
\draw[arrow] (linear)  -- (softmax);

% ── Nota máscara ─────────────────────────────────────────────
\node[draw=gray!30, fill=gray!5, rounded corners, font=\tiny,
      inner sep=4pt, align=left] at (6.8, 2.75)
    {\textbf{Masked:} m\'ascara causal\\impide ver tokens futuros.\\Garantiza generaci\'on\\autoregresiva.};
\draw[-{Stealth}, gray!40, thin] (4.9, 2.75) -- (5.5, 2.75);

\end{tikzpicture}
\caption{Arquitectura completa del Transformer original~\cite{vaswani2017attention}.
El encoder (izquierda) procesa la secuencia de entrada en paralelo y produce
representaciones contextualizadas. El decoder (derecha) genera la secuencia
de salida token a token, utilizando atenci\'on enmascarada sobre su propia
salida parcial y atenci\'on cruzada (\textit{cross-attention}) sobre las
representaciones del encoder. Las flechas rojas indican el flujo de $K$ y $V$
del encoder al decoder.}
\label{fig:transformer_completo}
\end{figure}
